%% appD --- Further reading
%%
%% Bibliographic details are language-independent and are identical in the
%% two editions; only the comments on them are translated. Every entry here
%% was checked against the work itself or its publisher's record; a work the
%% author could not check is not listed.
%%
%% Twin: appendices/pl/appD-reading.tex.

\chapter{Further reading}
\label{app:reading}

\noindent
Two lists: the papers in which the ideas of this book first appeared, and
the books to read next. The papers are for looking up where something came
from; most of them need more mathematics than this book assumes, but their
first pages are often clearer than anything written about them since. The
books are for going further, and each comes with a note on what it is good
for and what it asks of you.

\section{Where the ideas first appeared}
\label{sec:appd-papers}

\begin{description}[leftmargin=1.2em, itemsep=4pt, font=\normalfont\bfseries, before=\raggedright]
\item[Pierre-Simon Laplace,] \emph{Essai philosophique sur les
  probabilités}, 1814. The clockwork universe stated at its most confident;
  quoted in Chapter~\ref{ch:models}.
\item[Henri Poincaré,] \emph{Science et méthode}, 1908. The first clear
  statement, in words, that small differences at the start can produce
  very large ones in the end, and that prediction then becomes impossible.
\item[Edward N. Lorenz,] \enquote{Deterministic nonperiodic flow},
  \emph{Journal of the Atmospheric Sciences} 20 (1963), 130--141. The three
  equations of Chapter~\ref{ch:lorenz}, and the paper that started the modern
  subject.
\item[Tien-Yien Li and James A. Yorke,] \enquote{Period three implies chaos},
  \emph{American Mathematical Monthly} 82 (1975), 985--992. The paper that
  gave the subject its name (Chapter~\ref{ch:bifurcations}).
\item[Robert M. May,] \enquote{Simple mathematical models with very
  complicated dynamics}, \emph{Nature} 261 (1976), 459--467. The logistic
  map as a warning to everyone who builds models (Chapter~\ref{ch:logistic}).
\item[Michel Hénon,] \enquote{A two-dimensional mapping with a strange
  attractor}, \emph{Communications in Mathematical Physics} 50 (1976),
  69--77. The map of Chapter~\ref{ch:strange}.
\item[Mitchell J. Feigenbaum,] \enquote{Quantitative universality for a class
  of nonlinear transformations}, \emph{Journal of Statistical Physics} 19
  (1978), 25--52. The constant of Chapter~\ref{ch:bifurcations}.
\item[Benoit B. Mandelbrot,] \enquote{How long is the coast of Britain?
  Statistical self-similarity and fractional dimension}, \emph{Science} 156
  (1967), 636--638. Fractal dimension, before the word fractal
  (Chapter~\ref{ch:fractals}).
\item[Floris Takens,] \enquote{Detecting strange attractors in turbulence},
  in \emph{Dynamical Systems and Turbulence}, Lecture Notes in Mathematics
  898 (1981), 366--381. Rebuilding a system's shape from one measured
  quantity (Chapter~\ref{ch:life}).
\item[Jacques Laskar,] \enquote{A numerical experiment on the chaotic
  behaviour of the Solar System}, \emph{Nature} 338 (1989), 237--238. Chaos
  in the orbits of the planets (Chapter~\ref{ch:mechanics}).
\item[Edward Ott, Celso Grebogi and James A. Yorke,] \enquote{Controlling
  chaos}, \emph{Physical Review Letters} 64 (1990), 1196--1199.
  Chapter~\ref{ch:taming}.
\item[Louis M. Pecora and Thomas L. Carroll,] \enquote{Synchronization in
  chaotic systems}, \emph{Physical Review Letters} 64 (1990), 821--824.
  Chapter~\ref{ch:taming}.
\item[Edward N. Lorenz,] \enquote{Predictability: a problem partly solved},
  \emph{Seminar on Predictability}, European Centre for Medium-Range Weather
  Forecasts, 1996. The toy atmosphere of Chapter~\ref{ch:weather}.
\end{description}

\section{Books to read next}
\label{sec:appd-books}

\begin{description}[leftmargin=1.2em, itemsep=4pt, font=\normalfont\bfseries, before=\raggedright]
\item[James Gleick,] \emph{Chaos: Making a New Science}, 1987. The story of
  the people who found chaos, told as history rather than mathematics. The
  best book to read alongside this one.
\item[Edward N. Lorenz,] \emph{The Essence of Chaos}, 1993. One of the
  founders explaining the subject to a general reader, gently and with
  care; a natural next step after Parts~II and~III.
\item[Ian Stewart,] \emph{Does God Play Dice? The Mathematics of Chaos},
  1989. A lively popular account by a mathematician, with more mathematics
  than Gleick and less than a textbook.
\item[Leonard A. Smith,] \emph{Chaos: A Very Short Introduction}, 2007.
  Short and careful, and especially good on what chaos does and does not
  mean for forecasting.
\item[Tim Palmer,] \emph{The Primacy of Doubt}, 2022. Ensemble forecasting
  and the uses of uncertainty, from a meteorologist who helped build it;
  the book to read after Chapter~\ref{ch:weather}.
\item[Benoit B. Mandelbrot,] \emph{The Fractal Geometry of Nature}, 1982.
  The founding book of fractals, discursive and full of pictures.
\item[Heinz-Otto Peitgen, Hartmut Jürgens and Dietmar Saupe,] \emph{Chaos
  and Fractals: New Frontiers of Science}, 1992. A very large, very
  generous book that builds everything from experiments you can do; the
  natural continuation of this one for a reader who liked the labs.
\item[Steven H. Strogatz,] \emph{Nonlinear Dynamics and Chaos}, 1994. The
  textbook to read when you are ready for calculus: clear, applied and
  much loved. It is where this book's subject becomes a university course.
\item[Kathleen T. Alligood, Tim D. Sauer and James A. Yorke,] \emph{Chaos: An
  Introduction to Dynamical Systems}, 1996. A mathematics textbook written by
  one of the subject's founders and his colleagues, with the proofs.
\end{description}
